\section{Introduction}
Consider the reflection of a blast wave after it strikes a rigid target, Figure \ref{fig:RPB}. The angle of incidence $\theta$, at a point on the surface is defined as the angle between the outward normal of the surface and the direct vector from the explosive charge to that point. \cite{rigbyAngleIncidenceEffects2015}.
\begin{figure}[htpb]
	\centering
	\includegraphics[scale=1.5, keepaspectratio]{Images/RPB.pdf}
	\caption{Angle of incidence and slant distance for a specific point on a square target}
	\label{fig:RPB}
\end{figure}
When a blast wave interacts with a structure located in close proximity to it, the resulting blast load is extremely high-magnitude, short duration and highly spatially non-uniform. 
The reason for this spatial non-uniformity is due to charge properties (shape, type), mach stem effects and the varying standoff  distance comparative to the perpendicular scenario for the incident shock wave.

A significant challenge to the blast protection community is the quantification of loading from near-field explosive events, and modelling of the subsequent structural response.
In the near-field, loading is extremely sensitive to scaled distance and charge shape, fundamental properties of the explosive event. 

Physics based solvers can provide highly accurate solutions to specific scenarios but may take days to complete, conversely we have simplistic methods that allow risk-based design approaches, but these are not accurate and can only offer vague approximations.

This paper details a methodology for producing a FREM (fast running engineering model) for a given range of scaled distances that allows the blast protection community to perform risk-based engineering analysis. 
The produced models combine the accuracy of physics-based solutions (with experimental validation) and the speed of semi-empirical simplistic methods.